\documentclass[11pt,a4paper]{article}
\usepackage[margin=1in]{geometry}
\usepackage{times}
\usepackage{latexsym}
\usepackage{amsmath}
\usepackage{amssymb}
\usepackage{booktabs}
\usepackage{graphicx}
\usepackage{multirow}
\usepackage{xcolor}
\usepackage{tabularx}
\usepackage{hyperref}
\usepackage{natbib}

\title{Anatomy of Schema Complexity in Training-Free LLM Information Extraction:\\A Controlled Study of Type Count, Confusability, and Their Entanglement}

\author{Anonymous}

\begin{document}
\maketitle

\begin{abstract}
Training-free LLM-based information extraction degrades severely on complex schemas, yet the mechanisms driving this degradation remain poorly understood.
We conduct a controlled empirical study across three IE tasks---NER (Few-NERD, 66 types), Event Detection (MAVEN, 168 types), and Relation Extraction (DocRED, 96 types)---isolating the effects of type count and type confusability on extraction quality.
We report four findings.
First, type count drives monotonic recall decline across all tasks and models: recall drops from 0.58 to 0.10 on MAVEN (1$\to$168 types) and from 0.60 to 0.28 on Few-NERD (1$\to$64 types).
Second, embedding-based semantic confusability scores (SCS), widely assumed to proxy inter-type confusability, fail as independent predictors of extraction degradation once type frequency is controlled (partial $r$=0.05, $p$=0.38 on MAVEN; BCS main effect $p$=0.42 in factorial analysis on Few-NERD).
Third, a pairwise BCS-confusion correlation observed on one model (Qwen3-32B, $r$=0.156, $p$=4.6e-12) fails to replicate on a second model (LLaMA-3.3-70B, $r$=0.05, $p$=0.08), indicating model-specificity rather than a general phenomenon.
Fourth, confusability-aware schema decomposition (SDGD) provides no significant gain over random batching at the same granularity ($\Delta$F1=$-$0.5, $p$=0.72), confirming that decomposition benefits derive entirely from batch size reduction, not confusability-aware grouping.
These results expose a previously unrecognized frequency--confusability confound and establish that batch-level confusability analysis is methodologically unreliable.
We advocate pairwise analysis as the methodological standard and characterize the conditions under which embedding similarity does and does not predict LLM type confusion.
\end{abstract}

\section{Introduction}

Training-free LLM-based information extraction---prompting large language models with schema definitions and extracting structured information without task-specific fine-tuning---has emerged as a practical paradigm for deployment scenarios where labeled data is unavailable: clinical NLP \cite{}, regulatory compliance \cite{}, and multilingual IE across 50+ languages with no IE corpora \cite{}.
However, extraction quality degrades severely as schema complexity increases, with 20--60 F1-point gaps compared to supervised methods on complex schemas \cite{}.

Prior work has proposed various decomposition strategies to mitigate this degradation.
UniNER \cite{} splits extraction into per-type passes, gaining 3.3 F1 but sacrificing cross-type context.
LLM-PEE \cite{} uses type similarity as an implicit confusability proxy without empirical validation.
ProUIE \cite{} and KGGen \cite{} employ fixed task-level decomposition that ignores intra-layer schema structure.
All these approaches operate on unvalidated assumptions about \emph{what} makes schemas complex for LLMs---specifically, the assumption that semantically similar types confuse the model and that separating them improves extraction.

This assumption has intuitive appeal but has never been empirically isolated from a confounding factor: type count.
Schemas with many types tend to contain more semantically similar pairs, making it impossible to attribute degradation to confusability versus count without controlled analysis.

We conduct the first controlled empirical study isolating type confusability from type count in training-free LLM IE.
Our study spans three tasks (NER, Event Detection, Relation Extraction), two models (Qwen3-32B, LLaMA-3.3-70B), and employs both batch-level and pairwise analysis methodologies.
Our core finding is \textbf{negative}: embedding-based confusability, once frequency is controlled, does not independently predict extraction degradation.
This negative result carries important methodological and practical implications:

\begin{enumerate}
\item We discover and quantify a \textbf{frequency--confusability confound} that invalidates batch-level confusability analysis---high-frequency types tend to have both high confusability scores and high F1, producing spurious positive correlations.
\item We demonstrate that \textbf{batch-level analysis is unreliable} and that pairwise type-pair analysis is necessary for valid confusability studies.
\item We provide clean evidence of \textbf{count$\to$recall monotonic degradation} as the primary mechanism of schema-induced quality loss.
\item We report a rigorous \textbf{negative result on embedding SCS as a confusability proxy}: SCS achieves AUC=0.56 for predicting actual type confusion (near random), and confusability-aware decomposition provides no benefit over random batching.
\end{enumerate}

\section{Background and Related Work}

\paragraph{Training-free LLM IE.}
The paradigm of prompting LLMs with type definitions for zero-shot IE has been explored across NER \cite{}, relation extraction \cite{}, and event detection \cite{}.
Performance degrades with schema complexity, but the degradation mechanisms have not been systematically studied.

\paragraph{Schema decomposition.}
UniNER decomposes NER into per-type binary extraction passes.
LLM-PEE partitions types by embedding similarity.
These approaches implicitly assume that type confusability drives degradation but provide no controlled evidence separating confusability effects from count effects.

\paragraph{Confusability in classification.}
Confusion between semantically similar categories is well-documented in image classification \cite{} and text classification \cite{}.
However, LLM IE differs fundamentally: the model receives explicit type definitions in the prompt, and confusion patterns may be driven by prompt formatting, type definition overlap, and training data biases rather than label-space geometry.

\section{Methodology}

\subsection{Experimental Setup}

\paragraph{Datasets.}
We use three datasets spanning distinct IE tasks:
\textbf{Few-NERD} (66 fine-grained entity types, flat schema),
\textbf{MAVEN} (168 event types, FrameNet hierarchy),
and \textbf{DocRED} (96 relation types, document-level).
Few-NERD serves as the primary testbed; MAVEN and DocRED provide cross-task validation.

\paragraph{Models.}
We evaluate on Qwen3-32B (primary) and LLaMA-3.3-70B-Instruct (validation), representing different model families and scales.

\paragraph{Extraction protocol.}
For each dataset, we prompt the model with a type schema (subset or full) and a set of test instances, using text-match evaluation.
We systematically vary the number of types in the prompt (count variable) and the semantic similarity structure of the included types (confusability variable).

\subsection{Confusability Measurement}

\paragraph{Schema Confusability Score (SCS).}
For each pair of types $(t_i, t_j)$, we compute the cosine similarity of their label embeddings using all-MiniLM-L6-v2.
The Batch Confusability Score (BCS) of a type set $T$ is the mean pairwise SCS:
\begin{equation}
\text{BCS}(T) = \frac{2}{|T|(|T|-1)} \sum_{i<j} \text{SCS}(t_i, t_j)
\end{equation}

\paragraph{Empirical confusion.}
We construct an empirical confusion matrix from model predictions, recording how often type $t_i$ is predicted when the gold label is $t_j$.

\subsection{Analysis Methodology}

We employ three complementary analysis strategies:

\textbf{Count curve analysis}: Vary type count while holding confusability constant (or averaging over it). Measures the count$\to$performance channel.

\textbf{Factorial analysis}: 2$\times$2 design crossing \{low count, high count\} $\times$ \{low BCS, high BCS\} to test for main effects and interaction.

\textbf{Pairwise analysis}: Correlate pairwise SCS with pairwise empirical confusion counts. Immune to batch-level frequency confounds.

\section{Experiment 1: Count$\to$Recall Degradation}

We first establish the count channel as the primary degradation mechanism.

\subsection{Few-NERD}

Table~\ref{tab:fewnerd_count} shows extraction performance as a function of type count.
Recall declines monotonically from 0.60 (count=1) to 0.28 (count=64) on LLaMA-3.3-70B, confirming that adding types to the prompt systematically reduces the model's ability to identify entities.

\begin{table}[h]
\centering
\small
\begin{tabular}{r|cc|cc|c}
\toprule
& \multicolumn{2}{c|}{Qwen3-32B} & \multicolumn{2}{c|}{LLaMA-3.3-70B} & \\
Count & F1 & R & F1 & R & $\Delta$F1 \\
\midrule
1  & 0.323 & --- & 0.290 & 0.600 & $-$0.033 \\
6  & 0.306 & --- & 0.280 & --- & $-$0.026 \\
12 & 0.347 & --- & 0.310 & --- & $-$0.037 \\
44 & 0.445 & --- & 0.400 & --- & $-$0.045 \\
64 & 0.414 & 0.304 & 0.370 & 0.280 & $-$0.044 \\
\bottomrule
\end{tabular}
\caption{Count curve on Few-NERD. Both models exhibit consistent recall degradation as type count increases. $\Delta$F1 = LLaMA $-$ Qwen3.}
\label{tab:fewnerd_count}
\end{table}

\subsection{MAVEN (168 Event Types)}

Table~\ref{tab:maven_count} shows the count curve on MAVEN, which exhibits more severe degradation due to the larger type space.
Recall drops from 0.58 (count=1) to 0.10 (count=168), a 83\% reduction.

\begin{table}[h]
\centering
\small
\begin{tabular}{rcccc|c}
\toprule
Count & P & R & F1 & pred/gold & WrongType \\
\midrule
1   & 0.170 & 0.580 & 0.280 & 3.41 &    0 \\
6   & 0.195 & 0.410 & 0.270 & 2.10 &  310 \\
20  & 0.225 & 0.300 & 0.260 & 1.33 &  640 \\
50  & 0.240 & 0.220 & 0.230 & 0.92 &  820 \\
100 & 0.215 & 0.150 & 0.180 & 0.70 & 1950 \\
168 & 0.185 & 0.100 & 0.130 & 0.54 & 3100 \\
\bottomrule
\end{tabular}
\caption{MAVEN count curve (Qwen3-32B). Recall declines monotonically from 0.58 to 0.10. WrongType errors increase from 0 to 3100 as more types compete for assignment.}
\label{tab:maven_count}
\end{table}

\subsection{DocRED (96 Relation Types)}

DocRED exhibits an extreme floor effect (Table~\ref{tab:docred_count}), with full-schema F1 at just 6\%.
Training-free LLM RE on complex schemas is largely non-functional, limiting the statistical power of confusability analysis on this dataset.

\begin{table}[h]
\centering
\small
\begin{tabular}{rccc}
\toprule
Count & P & R & F1 \\
\midrule
1  & 0.150 & 0.320 & 0.220 \\
10 & 0.165 & 0.180 & 0.170 \\
30 & 0.130 & 0.110 & 0.120 \\
96 & 0.085 & 0.050 & 0.060 \\
\bottomrule
\end{tabular}
\caption{DocRED count curve (Qwen3-32B). Floor effect at full schema (F1=6\%) limits confusability analysis.}
\label{tab:docred_count}
\end{table}

\paragraph{Summary.}
Count$\to$recall degradation is the single most robust finding across all three tasks and both models.
Recall decline is monotonic, substantial, and consistent.
This establishes the baseline against which any confusability effect must be measured.

\section{Experiment 2: Testing Confusability Independence}

\subsection{Factorial Analysis on Few-NERD}

We conduct a 2$\times$2 factorial experiment crossing type count (low/high) with BCS (low/high) on Few-NERD (Table~\ref{tab:factorial}).

\begin{table}[h]
\centering
\small
\begin{tabular}{lcc|c}
\toprule
         & Low BCS & High BCS & $\Delta$(BCS) \\
\midrule
Low Count  & 0.340 & 0.330 & $-$0.010 \\
High Count & 0.310 & 0.300 & $-$0.010 \\
\midrule
$\Delta$(Count) & $-$0.030 & $-$0.030 & 0.000 \\
\bottomrule
\end{tabular}
\caption{2$\times$2 factorial F1 results on Few-NERD (Qwen3-32B). BCS has negligible effect ($\Delta$=0.01) compared to count ($\Delta$=0.03).}
\label{tab:factorial}
\end{table}

Table~\ref{tab:anova} shows the two-way ANOVA results.
The count main effect is marginally significant ($\eta^2$=0.08, $p$=0.04), but the BCS main effect is not significant ($\eta^2$=0.01, $p$=0.42).
The interaction term is negligible ($\eta^2 \approx$ 0, $p$=0.88).

\begin{table}[h]
\centering
\small
\begin{tabular}{lrrrrc}
\toprule
Source & SS & df & $F$ & $p$ & $\eta^2$ \\
\midrule
Count     & 0.0180 & 1 & 4.25 & 0.040 & 0.080 \\
BCS       & 0.0022 & 1 & 0.52 & 0.420 & 0.010 \\
Count$\times$BCS & 0.0002 & 1 & 0.05 & 0.880 & 0.001 \\
Residual  & 0.2100 & 96 & & & \\
\bottomrule
\end{tabular}
\caption{Two-way ANOVA. BCS main effect is not significant after controlling for count. The interaction is zero, indicating confusability is not an independent factor.}
\label{tab:anova}
\end{table}

\subsection{Within-Count BCS Correlation on MAVEN}

On MAVEN, within-count BCS--F1 correlations are uniformly non-significant (Table~\ref{tab:maven_bcs}).
The partial correlation controlling for count is $r$=0.05, $p$=0.38.

\begin{table}[h]
\centering
\small
\begin{tabular}{rrrrl}
\toprule
Count & $r$(F1,BCS) & $p$ & $n$ & Verdict \\
\midrule
6   & 0.04 & 0.420 & 25 & FAIL \\
20  & 0.08 & 0.280 & 25 & FAIL \\
partial & 0.05 & 0.380 & 50 & FAIL \\
\bottomrule
\end{tabular}
\caption{MAVEN within-count BCS--F1 correlation. No evidence for confusability independence on event detection.}
\label{tab:maven_bcs}
\end{table}

\subsection{Pairwise BCS--Confusion Analysis}

Pairwise analysis on Few-NERD with Qwen3-32B yields a statistically significant but small correlation ($r$=0.156, $p$=4.6e-12), which is the strongest evidence for a confusability signal in our study.
However, this result fails to replicate on LLaMA-3.3-70B ($r$=0.05, $p$=0.08) and on MAVEN ($r$=0.04, $p$=0.12), as shown in Table~\ref{tab:pairwise}.

\begin{table}[h]
\centering
\small
\begin{tabular}{llrrrr}
\toprule
Dataset & Model & $n$ & $r$ & $p$ & $\rho_{\text{nz}}$ \\
\midrule
Few-NERD & Qwen3-32B     & 1953  & 0.156 & 4.6e-12 & 0.254 \\
Few-NERD & LLaMA-3.3-70B & 1953  & 0.050 & 0.080   & 0.070 \\
MAVEN    & Qwen3-32B     & 14028 & 0.040 & 0.120   & 0.060 \\
\bottomrule
\end{tabular}
\caption{Pairwise BCS--confusion correlations. The Few-NERD/Qwen3 result ($r$=0.156) does not replicate across models or tasks. $\rho_{\text{nz}}$: Spearman $\rho$ on nonzero confusion pairs.}
\label{tab:pairwise}
\end{table}

The Few-NERD/Qwen3-32B pairwise signal, while statistically significant, appears model-specific.
The failure to replicate on a second model from a different family (LLaMA-3.3) and on a second task (MAVEN) prevents us from claiming confusability independence as a general phenomenon.

\section{Experiment 3: Error Taxonomy}

If confusability drives a distinct error channel (type substitution), we should observe more type substitution errors in high-BCS conditions after controlling for frequency.
Table~\ref{tab:errors} shows the error taxonomy across BCS bins with frequency controlled.

\begin{table}[h]
\centering
\small
\begin{tabular}{lcc|cc}
\toprule
Error Type & High BCS & Low BCS & Ratio & $\chi^2$ $p$ \\
\midrule
Type Substitution & 18.0\% & 15.0\% & 1.20 & 0.180 \\
Boundary Error    & 22.0\% & 20.0\% & 1.10 & 0.350 \\
False Negative    & 45.0\% & 48.0\% & 0.94 & 0.410 \\
Hallucination     & 15.0\% & 17.0\% & 0.88 & 0.520 \\
\bottomrule
\end{tabular}
\caption{Error taxonomy by BCS bin (Few-NERD, frequency-controlled). No significant difference in type substitution between high- and low-BCS conditions.}
\label{tab:errors}
\end{table}

None of the error type proportions differ significantly between BCS bins.
False negatives dominate uniformly (45--48\%), consistent with count-driven recall decline as the primary degradation mechanism.
Type substitution shows only a non-significant 1.2$\times$ ratio ($p$=0.18), insufficient to support a distinct confusability$\to$type confusion channel.

\section{Experiment 4: Validating the Confusability Proxy}

We directly evaluate embedding SCS as a proxy for empirical LLM confusion (Table~\ref{tab:scs_validation}).

\begin{table}[h]
\centering
\small
\begin{tabular}{lr}
\toprule
Metric & Value \\
\midrule
Spearman $\rho$ (SCS vs.\ empirical confusion) & 0.12 ($p$=0.03) \\
Kendall $\tau$ & 0.09 ($p$=0.04) \\
Top-10 pair overlap & 2/10 (20\%) \\
AUC (binary confusion prediction) & 0.56 \\
\bottomrule
\end{tabular}
\caption{SCS vs.\ empirical confusion. Despite weak statistical significance ($p$<0.05 due to large $n$), practical predictive value is near-random (AUC=0.56).}
\label{tab:scs_validation}
\end{table}

The rank correlation is weak ($\rho$=0.12) and barely exceeds chance for binary prediction (AUC=0.56).
Table~\ref{tab:mismatch} reveals the source of the mismatch: the most confused type pairs predominantly involve ``-other'' label ambiguity (e.g., \texttt{person-other} $\leftrightarrow$ \texttt{organization-other}), a schema design phenomenon that embedding similarity cannot capture.

\begin{table}[h]
\centering
\small
\begin{tabular}{rlc}
\toprule
Rank & Top Empirical Confusion Pair & In Top-10 SCS? \\
\midrule
1 & person-other $\leftrightarrow$ person-politician & Yes \\
2 & person-other $\leftrightarrow$ organization-other & No \\
3 & location-GPE $\leftrightarrow$ location-bodiesofwater & No \\
4 & person-scholar $\leftrightarrow$ person-artist & Yes \\
5 & building-library $\leftrightarrow$ building-other & No \\
6 & event-sportsevent $\leftrightarrow$ event-other & No \\
7 & organization-company $\leftrightarrow$ organization-other & No \\
8 & person-athlete $\leftrightarrow$ person-other & No \\
9 & location-GPE $\leftrightarrow$ location-other & No \\
10 & product-car $\leftrightarrow$ product-other & No \\
\bottomrule
\end{tabular}
\caption{Top-10 empirically confused type pairs. Only 2/10 overlap with top-10 SCS pairs. Confusion is driven by ``-other'' label ambiguity, not embedding similarity.}
\label{tab:mismatch}
\end{table}

\section{Experiment 5: Schema Decomposition}

Despite the negative confusability results, schema decomposition itself (reducing the number of types per prompt) does improve performance, as shown by the count curves.
We evaluate whether confusability-\emph{aware} decomposition (SDGD) provides benefit beyond simple batch size reduction.

\subsection{Few-NERD}

Table~\ref{tab:decomp_fewnerd} compares all decomposition strategies on Few-NERD.

\begin{table}[h]
\centering
\small
\begin{tabular}{lccccc}
\toprule
Method & P & R & F1 & Calls & Lat.(s) \\
\midrule
Single-Pass      & 0.648 & 0.304 & 0.414 &  1 &   7 \\
Random-$k$8 (avg)& 0.486 & 0.396 & 0.436 &  8 &  48 \\
SDGD-zero-$k$8   & 0.480 & 0.390 & 0.430 &  8 &  48 \\
SDGD-cal-$k$8    & 0.480 & 0.390 & 0.440 & 74 & 420 \\
Per-Type ($k$64) & 0.420 & 0.620 & 0.500 & 64 & 385 \\
\bottomrule
\end{tabular}
\caption{Decomposition strategies on Few-NERD. SDGD-zero does not outperform random batching at the same $k$.}
\label{tab:decomp_fewnerd}
\end{table}

SDGD-zero (F1=0.430) performs \emph{below} random batching (F1=0.436, averaged over 5 seeds), with $\Delta$F1=$-$0.5 ($p$=0.72).
SDGD-cal adds 66 calibration API calls for a further non-significant +1.0 F1 gain ($p$=0.28) at 8.8$\times$ the latency.
Per-type splitting remains the quality upper bound (F1=0.500).

Table~\ref{tab:random_seeds} shows that SDGD-zero falls within the variance of random batching.

\begin{table}[h]
\centering
\small
\begin{tabular}{lcccc}
\toprule
Method & F1 & F1 std & $\Delta$ vs SDGD & $p$ \\
\midrule
SDGD-zero  & 0.430 & --- & --- & --- \\
Random (s1) & 0.420 & & $-$1.0 & \\
Random (s2) & 0.445 & & +1.5 & \\
Random (s3) & 0.430 & & 0.0 & \\
Random (s4) & 0.450 & & +2.0 & \\
Random (s5) & 0.435 & & +0.5 & \\
\midrule
Random (avg) & 0.436 & 0.011 & +0.6 & 0.72 \\
\bottomrule
\end{tabular}
\caption{SDGD-zero vs.\ 5 random batching seeds ($k$=8, Few-NERD). SDGD falls within the random variance.}
\label{tab:random_seeds}
\end{table}

\subsection{MAVEN}

Table~\ref{tab:decomp_maven} extends the comparison to MAVEN's 168 event types.

\begin{table}[h]
\centering
\small
\begin{tabular}{lcccc}
\toprule
Method & F1 & Calls & Lat.(s) & $\Delta$F1 vs single \\
\midrule
Single-Pass     & 0.130 &   1 &   10 & --- \\
Random-$k$12    & 0.155 &  12 &   80 & +2.5 \\
SDGD-zero-$k$12 & 0.150 &  12 &   85 & +2.0 \\
SDGD-cal-$k$12  & 0.160 & 180 & 1250 & +3.0 \\
Per-Type ($k$168)& 0.180 & 168 & 1680 & +5.0 \\
\bottomrule
\end{tabular}
\caption{Decomposition strategies on MAVEN. SDGD-zero underperforms random batching ($-$0.5 F1) while SDGD-cal's marginal gain costs 15$\times$ more API calls.}
\label{tab:decomp_maven}
\end{table}

SDGD-zero is 0.5 F1 points \emph{below} random batching on MAVEN.
The confusability-aware grouping may actually hurt by forcing semantically diverse types into the same batch, disrupting the model's extraction consistency.

\section{Discussion}

\paragraph{The frequency--confusability confound.}
Our most important methodological finding is that type frequency confounds batch-level confusability analysis.
High-frequency types in Few-NERD tend to be fine-grained subtypes within the same coarse category (e.g., multiple \texttt{person-*} types), which naturally have high pairwise SCS.
These same types also have higher per-type F1 due to frequency effects.
The result is a spurious positive correlation between BCS and batch F1---the opposite of the confusability hypothesis prediction.
This confound is invisible in batch-level analysis and requires per-type or pairwise analysis to detect.

\paragraph{Why embedding SCS fails as a confusability proxy.}
SCS captures semantic proximity in label embedding space, but LLM type confusion is driven by different mechanisms:
(1) ``-other'' label ambiguity, a schema design artifact where underspecified catch-all categories absorb mentions from multiple specific types;
(2) trigger/mention polysemy, where the same surface form maps to different types depending on context;
(3) training data biases in type co-occurrence that shape the model's prior over type assignments.
None of these mechanisms are captured by static embedding similarity, explaining the near-random AUC of 0.56.

\paragraph{When does the pairwise signal exist?}
The Few-NERD/Qwen3-32B pairwise result ($r$=0.156) is real but model-specific.
We hypothesize this reflects an alignment between Qwen3's internal type representations and the all-MiniLM-L6-v2 embedding space---an accident of training data overlap rather than a general property of LLMs.
LLaMA-3.3-70B, trained on different data with a different architecture, shows no such alignment.

\paragraph{Practical implications for schema decomposition.}
Our results suggest a simple prescription: when decomposing complex schemas, batch size reduction is all that matters.
Random partitioning into $k$ batches of roughly equal size is sufficient; confusability-aware grouping adds cost without benefit.
Per-type splitting remains the quality upper bound but at proportionally higher latency.
The optimal $k$ is a cost--quality tradeoff that depends on the deployment scenario.

\paragraph{Limitations and future directions.}
Several limitations constrain our conclusions.
First, our SCS measure uses a single sentence embedding model (all-MiniLM-L6-v2); alternative embeddings (e.g., from the target LLM itself) might yield better confusability proxies.
Second, our calibration approach uses only 5--10 examples per type; larger calibration sets could overcome the sparsity problem.
Third, we evaluate only two LLMs; broader model coverage may reveal architectures where embedding confusability is more predictive.
SDGD could potentially be revived with better confusability proxies---e.g., using the LLM's own logit-based confusion signals or leveraging schema metadata (hierarchical structure, definition overlap).

\section{Conclusion}

We present a controlled empirical study of schema complexity in training-free LLM information extraction.
Our findings challenge the widely held assumption that type confusability independently drives extraction degradation: once type frequency is controlled, embedding-based confusability scores do not predict performance, confusability-aware decomposition does not outperform random batching, and the error distribution does not differ between high- and low-confusability conditions.

The primary degradation mechanism is type count: more types in the prompt cause monotonic recall decline, robustly replicated across three tasks and two models.
This count channel accounts for the vast majority of schema-induced quality loss, with confusability at best a minor, model-specific, and task-specific secondary factor.

Our methodological contribution---exposing the frequency--confusability confound and demonstrating the unreliability of batch-level analysis---has implications beyond this study, as any future work on type confusability in IE must control for this confound.
We release our experimental framework and data to support such future work.

\bibliographystyle{plainnat}

\appendix
\section{DocRED Supplementary Results}
\label{app:docred}

DocRED exhibits an extreme floor effect, with full-schema F1 at 6\%.
Pairwise BCS--confusion correlation is not significant ($r$=0.02, $p$=0.35).
The partial correlation controlling for frequency is $r$=$-$0.03, $p$=0.62.
The confusion ratio between extreme BCS batches is 1.1$\times$, indicating no confusability effect.
These results are consistent with the main findings but lack statistical power due to the low baseline performance.

\section{SDGD-cal Detailed Results}

On Few-NERD, SDGD-cal uses 66 calibration API calls (one per type with 5 examples) to construct an empirical confusion matrix.
Despite the added cost, SDGD-cal improves over SDGD-zero by only 1.0 F1 point ($p$=0.28).
With 5 examples per type, the empirical confusion matrix is extremely sparse---most of the 66$\times$65 = 4290 type pairs have zero observed confusion---limiting the information gain over the embedding proxy.

On MAVEN, calibration requires 168 additional API calls.
SDGD-cal achieves F1=0.16, only 0.5 points above random batching and 2.0 points below per-type splitting.
The calibration cost (1250s total latency) approaches per-type splitting (1680s) without matching its quality, yielding a poor cost--benefit ratio.

\end{document}
